\documentclass[11pt]{article}
\usepackage[a4paper,margin=26mm]{geometry}
\usepackage{amsmath,amssymb,mathtools}
\usepackage{unicode-math}
\setmainfont{Latin Modern Roman}
\setmathfont{Latin Modern Math}
\usepackage{xurl,hyperref}
\hypersetup{hidelinks,pdftitle={Problem statement}}
\newcommand{\sref}[2]{\hyperlink{source:#1:#2}{[S#2]}}
\newcommand{\cataloguescope}[1]{#1}
\begin{document}
\section{Classical 3-manifold theory — Question 9.3}\label{statement}

\subsection{Definitions and mathematical statement}

\paragraph{Shared notation.}$\mathbb N=\{1,2,\ldots\}$, and $\mathbb Z,\mathbb Q,\mathbb R,\mathbb C$ denote the integers, rationals, reals and complex numbers. Any other conventions are specified locally.


A surface lamination $\Lambda$ is a closed subset of $M$ locally homeomorphic to a product of a disk with a closed transverse set. It is essential when it has no sphere leaf or torus leaf bounding a solid torus, and its complementary regions are irreducible, have incompressible boundary, and have no compressing monogon. Such a monogon is a compressing disk with one boundary arc on a boundary leaf and the other along a transverse interval fiber. It is genuine if some complementary region is not an interval bundle. A sublamination is a closed union of leaves. A persistent lamination in a knot exterior remains essential after every nontrivial Dehn filling under consideration; meridional filling is the trivial surgery. Delman--Roberts methods construct such laminations from knot projections, especially alternating projections.
\paragraph{Statement recorded in the supplied source.}
Determine the largest classes of knots for which the Delman--Roberts projection-based techniques can be extended to construct persistent essential laminations. The task concerns the scope of these constructions and their effective extension, not merely existence of an arbitrary essential lamination in a knot complement. \sref{statement}{2}
\subsection{Short English statement}
Extend the projection-based Delman--Roberts construction of persistent laminations to the widest possible classes of knots.
\subsection{Sources}
\begin{description}
\item[\hypertarget{source:statement:1}{S1}] ULAM.AI and the UnsolvedMath Contributors, \emph{UnsolvedMath: A Curated Collection of Open Mathematics Problems}, catalogue entry. CC BY 4.0. \url{https://huggingface.co/datasets/ulamai/UnsolvedMath}.
\item[\hypertarget{source:statement:2}{S2}] Danny Calegari, \emph{Problems in foliations and laminations of 3-manifolds}, version 0.78 (2002), Question 9.3; Definitions 1.1--1.10 and local definitions and remarks. \url{https://arxiv.org/abs/math/0209081v1}.
\end{description}
\label{end:statement}
\end{document}
